\documentclass[11pt,a4paper]{article}

\usepackage[utf8]{inputenc}
\usepackage[T1]{fontenc}
\usepackage[spanish,es-noindentfirst]{babel}
\usepackage[margin=2.5cm]{geometry}
\usepackage{graphicx}
\usepackage{xcolor}
\usepackage{listings}
\usepackage{tikz}
\usepackage{hyperref}
\usepackage{fancyhdr}
\usepackage{enumitem}
\usepackage{longtable}
\usepackage{booktabs}
\usepackage{array}
\usetikzlibrary{arrows.meta, positioning, calc, shapes.geometric}

\definecolor{codebg}{RGB}{245,245,245}
\definecolor{codekeyword}{RGB}{0,0,150}
\definecolor{codestring}{RGB}{140,0,0}
\definecolor{codecomment}{RGB}{0,120,0}
\definecolor{uteqverde}{RGB}{0,90,60}
\definecolor{alerta}{RGB}{150,30,30}
\definecolor{okverde}{RGB}{0,110,60}

\lstdefinestyle{codigo}{
  backgroundcolor=\color{codebg},
  basicstyle=\ttfamily\footnotesize,
  keywordstyle=\color{codekeyword}\bfseries,
  stringstyle=\color{codestring},
  commentstyle=\color{codecomment}\itshape,
  numbers=left,
  numberstyle=\tiny\color{gray},
  breaklines=true,
  frame=single,
  rulecolor=\color{black!30},
  tabsize=2,
  showstringspaces=false,
  captionpos=b
}
\lstset{style=codigo}

\hypersetup{
  colorlinks=true,
  linkcolor=uteqverde,
  urlcolor=uteqverde,
  citecolor=uteqverde,
  pdftitle={SGED - Informe de la Tercera Entrega},
  pdfauthor={Arcalle Grefa Darwin Orlando, Pallo Pinto Alejandro Daniel, Velez Lopez Ricardo Elias}
}

\pagestyle{fancy}
\fancyhf{}
\setlength{\headheight}{14pt}
\renewcommand{\headrulewidth}{0.4pt}
\lhead{SGED --- Tercera Entrega}
\rhead{Aplicaciones Web --- PPA 2026-2027}
\cfoot{\thepage}

\newcommand{\ruta}[1]{\texttt{\footnotesize #1}}
\newcommand{\ok}{\textcolor{okverde}{\textbf{OK}}}
\newcommand{\pendiente}{\textcolor{alerta}{\textbf{pendiente}}}

\begin{document}

% ============================================================
% PORTADA
% ============================================================
\begin{titlepage}
\centering
\vspace*{1cm}

{\LARGE\bfseries UNIVERSIDAD TÉCNICA ESTATAL DE QUEVEDO}\\[0.3cm]
{\large Facultad de Ciencias de la Ingeniería}\\[0.15cm]
{\large Carrera de Ingeniería en Software}\\[2cm]

{\Huge\bfseries\color{uteqverde} SGED}\\[0.4cm]
{\Large Sistema de Gestión para la\\Escuela Deportiva ProFútbol}\\[1.5cm]

{\LARGE\bfseries Informe de la Tercera Entrega}\\[0.4cm]
{\large Proyecto Fin de Curso --- Aplicaciones Web}\\[2cm]

\begin{tabular}{rl}
\textbf{Asignatura:} & Aplicaciones Web (Quinto nivel) \\[0.2cm]
\textbf{Docente:} & Dr. Gleiston Cicerón Guerrero Ulloa, Ph.D. \\[0.2cm]
\textbf{Periodo:} & PPA 2026-2027 \\[0.2cm]
\textbf{Versión:} & \texttt{v0.9.0-rc} (candidata) \\
\end{tabular}\\[1.5cm]

{\large\bfseries Integrantes:}\\[0.3cm]
{\large Arcalle Grefa Darwin Orlando}\\
{\large Pallo Pinto Alejandro Daniel}\\
{\large Velez Lopez Ricardo Elias}\\[1.5cm]

{\large\bfseries Repositorio:}\\[0.2cm]
{\ttfamily\url{https://github.com/DarwinSM21/SGED_APPWEB}}\\[1.5cm]

{\large Quevedo --- Los Ríos --- Ecuador --- 2026}

\vfill
\end{titlepage}

\tableofcontents
\newpage

% ============================================================
\section{Introducción}

Este informe documenta la Tercera Entrega del Proyecto Fin de Curso de la
asignatura Aplicaciones Web: el sistema \textbf{SGED}, una aplicación web
para la gestión administrativa y deportiva de la escuela de fútbol formativo
ProFútbol.

El eje de esta entrega es la \emph{verificabilidad}. Todo lo que se afirma
aquí puede comprobarse ejecutando el repositorio: cada número proviene de una
medición archivada en \ruta{docs/mediciones/}, y cada requisito enlaza con el
archivo que lo implementa y con la prueba que lo cubre.

\subsection{Estado de la entrega: declaración previa}

Antes de exponer los resultados conviene ser explícito sobre dos cosas:

\begin{enumerate}[leftmargin=*]
  \item \textbf{Alcance realmente implementado.} El sistema expone hoy 13
    \emph{endpoints} REST (autenticación y gestión de estudiantes). El
    dominio deportivo (entrenadores, horarios, sesiones, asistencias y
    evaluaciones) está \emph{modelado y migrado} en la base de datos, pero
    no expuesto vía API. En el documento de requisitos cada elemento indica
    su estado real, y en este informe se mantiene la misma distinción. No se
    presenta como entregado nada que no lo esté.
  \item \textbf{Medición de usabilidad pendiente.} La encuesta SUS exigida
    por el Bloque~C.3 requiere al menos diez participantes externos. El
    instrumento y el script de análisis están construidos y versionados,
    pero \textbf{las respuestas todavía no han sido recolectadas}. No se
    reportan resultados inventados (sección~\ref{sec:sus}).
\end{enumerate}

\subsection{Pila tecnológica}

\begin{center}
\begin{tabular}{@{}ll@{}}
\toprule
\textbf{Capa} & \textbf{Tecnología} \\
\midrule
Backend & Spring Boot 3.2 (Java 21 LTS), Spring Data JPA, Spring Security \\
Frontend & Angular, servido por nginx con terminación TLS \\
Base de datos & PostgreSQL 16 (esquemas \texttt{seguridad} y \texttt{deportivo}) \\
Caché y revocación & Redis 7 \\
Migraciones & Flyway (\texttt{V1}--\texttt{V6}) \\
Orquestación & Docker Compose, imágenes fijadas por digest SHA-256 \\
\bottomrule
\end{tabular}
\end{center}

% ============================================================
\section{Resolución de las observaciones previas (Bloque 0)}

El docente emitió doce observaciones sobre las entregas 1A y 1B. Esta sección
documenta el estado de cada una con el \emph{commit} que la resuelve, de modo
que la afirmación sea verificable con \texttt{git show <hash>} y no
únicamente declarativa.

\renewcommand{\arraystretch}{1.25}
\begin{longtable}{@{}p{1.3cm}p{6.6cm}p{5.3cm}p{1.6cm}@{}}
\toprule
\textbf{Cód.} & \textbf{Observación del docente} & \textbf{Resolución verificable} & \textbf{Commit} \\
\midrule
\endfirsthead
\toprule
\textbf{Cód.} & \textbf{Observación} & \textbf{Resolución} & \textbf{Commit} \\
\midrule
\endhead

OBS-01 & Los requisitos se redactan como títulos («Registro de
estudiantes»); no usan «El sistema deberá\ldots». &
\ruta{docs/requisitos/SRS.md} reescribe los 22 RF y 13 RNF con la forma
normativa y criterio de verificación~\cite{iso29148}. & \texttt{98bab0d} \\

OBS-02 & Los diagramas C4 N1/N2 y el MER contienen texto marcador
«[Reemplazar con PNG]». & \ruta{docs/diagramas/} contiene los PNG reales
más las fuentes \texttt{.puml} y \texttt{.dbml}. & \texttt{bf6ee80} \\

OBS-03 & Incoherencia de pila: el ADR decide PHP/Laravel/MySQL, pero el
repositorio construye Spring Boot/Angular/PostgreSQL. & ADR-001 reescrito:
evalúa las tres opciones y decide explícitamente Java 21 con Spring Boot
3.2.5, coherente con el código. & \texttt{bf6ee80} \\

OBS-04 & \texttt{query.sql} con 0 claves foráneas y 88 columnas
«(PK)/(FK)» literales; exportación cruda no funcional. & \ruta{db/schema.sql}
declara 21 restricciones de integridad referencial reales. &
\texttt{db77ce2} \\

OBS-05 & Roles sin completar («Integrante 1/2/3»); wireframes con marca
ajena; falta \texttt{.env.example}. & \ruta{CONTRIBUTORS.md} asigna roles
CRediT por evidencia de \texttt{git log}; \ruta{.env.example} versionado. &
\texttt{a8d6f18} \\

OBS-06 & En el repositorio solo consta ADR-003; los diagramas aparecen en
el informe como texto. & Diagramas versionados como archivos; el
diccionario de datos completo se agrega en
\ruta{docs/basedatos/DATA-DICTIONARY.md}. & \texttt{98bab0d} \\

OBS-07 & \texttt{AuthController} solo expone \texttt{/login},
\texttt{/me} y \texttt{/ping}: no hay registro, logout ni refresh.
\texttt{RedisBlacklistService} existe pero no se invoca. & Seis
\emph{endpoints} operativos; el logout invoca realmente la lista de
revocación por \texttt{jti}. & \texttt{b907d10}, \texttt{bcb3642} \\

OBS-08 & La migración \texttt{V1} está vacía (0 bytes) y \texttt{V2}
depende de un esquema que ninguna migración crea; con
\texttt{ddl-auto=validate} el arranque no es reproducible. & \texttt{V1}
crea ambos esquemas y las tablas de identidad; el arranque desde clonación
limpia se verificó. & \texttt{b907d10} \\

OBS-09 & La lista de revocación no está cableada a ningún
\emph{endpoint} y no se aplica \texttt{@PreAuthorize}. & Ocho anotaciones
\texttt{@PreAuthorize} activas; revocación verificada por prueba
automatizada. & \texttt{688c4be}, \texttt{90c6c57} \\

OBS-10 & \texttt{AuthServiceTest} está vacío (0 bytes); la única prueba
real es \texttt{contextLoads()}. No se cumple el mínimo de 5 pruebas. La
colección Postman no está versionada. & 34 pruebas en 7 clases; cobertura
real 68,1\,\%; \ruta{docs/postman/coleccion.json} versionada. &
\texttt{1e798bc} \\

OBS-11 & \texttt{docker-compose.yml} está vacío (0 bytes); no hay
orquestación verificable. & Cuatro servicios orquestados con
\emph{healthchecks} e imágenes fijadas por digest; \texttt{make up} levanta
todo. & \texttt{b907d10}, \texttt{17dc5df} \\

OBS-12 & Varios contenidos del informe no se corresponden con lo presente
en el repositorio (están vacíos o ausentes). & Cada requisito declara su
estado real (implementado / modelado / planificado); esta entrega no afirma
nada no verificable. & \texttt{98bab0d} \\

\bottomrule
\end{longtable}
\renewcommand{\arraystretch}{1}

Las doce observaciones tienen resolución trazable. La más costosa de
resolver no fue una funcionalidad ausente sino OBS-12, porque exigió revisar
sistemáticamente qué afirmaba el informe anterior frente a lo que el
repositorio realmente contenía.

% ============================================================
\section{Arquitectura}

\subsection{Vista de contenedores}

\begin{center}
\begin{tikzpicture}[
  node distance=1.6cm,
  caja/.style={draw, rounded corners, fill=codebg, minimum width=3.1cm,
               minimum height=1.15cm, align=center, font=\small},
  flecha/.style={-{Latex[length=2.5mm]}, thick, color=uteqverde}
]
\node[caja] (nav) {Navegador\\\footnotesize (SPA Angular)};
\node[caja, right=2.2cm of nav] (ngx) {nginx\\\footnotesize :4200 / :8443 TLS};
\node[caja, right=2.2cm of ngx] (api) {Spring Boot\\\footnotesize :8080};
\node[caja, below=1.5cm of api] (pg) {PostgreSQL 16\\\footnotesize esquemas seguridad/deportivo};
\node[caja, left=2.2cm of pg] (rd) {Redis 7\\\footnotesize caché + revocación JTI};

\draw[flecha] (nav) -- node[above, font=\scriptsize]{HTTPS} (ngx);
\draw[flecha] (ngx) -- node[above, font=\scriptsize]{HTTP} (api);
\draw[flecha] (api) -- node[right, font=\scriptsize]{JPA} (pg);
\draw[flecha] (api) -- node[above, font=\scriptsize, sloped]{cookie JTI} (rd);
\end{tikzpicture}
\end{center}

La documentación de arquitectura en el modelo C4~\cite{brown2018c4} y las
decisiones arquitectónicas (ADR)~\cite{nygard2011adr} se mantienen en
\ruta{docs/diagramas/} y \ruta{docs/adr/}.

\subsection{Autenticación: JWT en cookie \texttt{HttpOnly}}

El sistema emite el JWT~\cite{rfc7519} exclusivamente en una cookie
\texttt{HttpOnly}, \texttt{Secure} y \texttt{SameSite=Strict}. El token nunca
viaja en el cuerpo de la respuesta ni se almacena en \texttt{localStorage},
de modo que un ataque de \emph{scripting} entre sitios no puede leerlo.

\begin{lstlisting}[caption={Emisión de la cookie de sesión (\ruta{AuthController.java})}]
ResponseCookie.from("sged_access", token)
        .httpOnly(true)
        .secure(cookieSecure)
        .sameSite("Strict")
        .path("/")
        .maxAge(Duration.ofMillis(expiracionMs))
        .build();
\end{lstlisting}

Como JWT es \emph{stateless}, cerrar sesión no invalidaría el token por sí
solo. Por eso el sistema mantiene una lista de revocación en Redis indexada
por el identificador del token (\texttt{jti}), con tiempo de vida igual al
tiempo que le restaba al token: un token cerrado deja de ser aceptado
inmediatamente, aunque su firma siga siendo válida.

% ============================================================
\section{Estrategia híbrida de acceso a datos}

El Bloque~A.2 de la guía exige una separación explícita entre lo que resuelve
el ORM y lo que debe resolver el motor de base de datos. La regla adoptada
es:

\begin{center}
\begin{tabular}{@{}p{7.3cm}p{7.3cm}@{}}
\toprule
\textbf{Spring Data JPA (ORM)} & \textbf{Procedimientos almacenados} \\
\midrule
CRUD elemental, consultas paginadas, búsquedas por identificador. &
Agregaciones, actualizaciones masivas con criterio de negocio, operaciones
multi-tabla transaccionales. \\
\bottomrule
\end{tabular}
\end{center}

Los dos procedimientos implementados están versionados en
\ruta{db/procs/} y creados por migración Flyway~\cite{flyway}:

\begin{center}
\begin{tabular}{@{}lll@{}}
\toprule
\textbf{Procedimiento} & \textbf{Requisito} & \textbf{Motivo} \\
\midrule
\texttt{sp\_contar\_estudiantes\_activos} & RF-14 & Agregación \texttt{COUNT} \\
\texttt{sp\_desactivar\_estudiantes\_categoria} & RF-15 & Actualización masiva transaccional \\
\bottomrule
\end{tabular}
\end{center}

\subsection{Un defecto real: \texttt{FUNCTION} frente a \texttt{PROCEDURE}}

La migración \texttt{V5} creó estos objetos como \texttt{FUNCTION}. Al
invocarlos desde Java con \texttt{@Procedure}, Spring Data usa
\texttt{CallableStatement} con sintaxis \texttt{\{call ...\}}, y PostgreSQL
solo acepta \texttt{CALL} contra procedimientos reales~\cite{postgresql16};
el error era explícito: \emph{«... is not a procedure. Hint: To call a
function, use SELECT.»}

La migración \texttt{V6} los convierte en \texttt{PROCEDURE} con parámetro de
salida. Se documenta porque ilustra la diferencia entre «el procedimiento
existe» y «el procedimiento se invoca de verdad»: en la versión anterior el
objeto SQL estaba creado pero quedaba huérfano, y la aplicación seguía
contando en memoria.

\begin{lstlisting}[caption={Invocación real desde el repositorio (\ruta{EstudianteRepository.java})}]
@Procedure(procedureName = "seguridad.sp_contar_estudiantes_activos")
Long contarActivosPorCategoria(@Param("p_categoria") String categoria);
\end{lstlisting}

\subsection{Ausencia de SQL dinámico}

Ninguna capa construye sentencias SQL por concatenación de cadenas. Toda
consulta usa parámetros vinculados (JPA) o parámetros nombrados (los
procedimientos). El repositorio incluye una auditoría automatizada
(\ruta{scripts/audit-sql-dynamic.sh}) integrada en \texttt{make audit}.

% ============================================================
\section{Requisitos y modelo de datos}

\subsection{Resumen de requisitos}

La especificación completa está en \ruta{docs/requisitos/SRS.md}. Se resumen
aquí los requisitos funcionales con su estado real y el \emph{endpoint} o
esquema que los materializa.

\renewcommand{\arraystretch}{1.15}
\begin{longtable}{@{}p{1.3cm}p{7.4cm}p{4.6cm}p{1.7cm}@{}}
\toprule
\textbf{ID} & \textbf{Requisito} & \textbf{Implementación} & \textbf{Estado} \\
\midrule
\endfirsthead
\toprule
\textbf{ID} & \textbf{Requisito} & \textbf{Implementación} & \textbf{Estado} \\
\midrule
\endhead

RF-01 & Registro de usuarios por administrador & \texttt{POST /api/auth/registro} & \ok \\
RF-02 & Autenticación con JWT en cookie \texttt{HttpOnly} & \texttt{POST /api/auth/login} & \ok \\
RF-03 & Cierre de sesión con revocación efectiva por \texttt{jti} & \texttt{POST /api/auth/logout} & \ok \\
RF-04 & Renovación de sesión sin re-credenciales & \texttt{POST /api/auth/refresh} & \ok \\
RF-05 & Consulta de la sesión activa & \texttt{GET /api/auth/me} & \ok \\
RF-06 & Bloqueo tras 5 intentos fallidos en 15 minutos & \texttt{LoginAttemptService} & \ok \\
RF-07 & Comprobación de disponibilidad & \texttt{GET /api/auth/ping} & \ok \\
RF-08 & Listado paginado de estudiantes & \texttt{GET /api/estudiantes} & \ok \\
RF-09 & Consulta por identificador con \texttt{404} tipificado & \texttt{GET /api/estudiantes/\{id\}} & \ok \\
RF-10 & Registro transaccional de estudiante & \texttt{POST /api/estudiantes} & \ok \\
RF-11 & Validación del formato de categoría (\texttt{SUB-NN}) & \texttt{EstudianteRequest} & \ok \\
RF-12 & Actualización de estudiante & \texttt{PUT /api/estudiantes/\{id\}} & \ok \\
RF-13 & Baja lógica preservando el historial & \texttt{DELETE /api/estudiantes/\{id\}} & \ok \\
RF-14 & Conteo de activos por categoría (en el motor) & \texttt{sp\_contar\_estudiantes\_activos} & \ok \\
RF-15 & Desactivación masiva por categoría (en el motor) & \texttt{sp\_desactivar\_estudiantes\_categoria} & \ok \\
\midrule
RF-16 & Gestión de entrenadores & \texttt{deportivo.entrenadores} & Modelado \\
RF-17 & Horarios recurrentes de entrenamiento & \texttt{horarios\_entrenamiento} & Modelado \\
RF-18 & Sesiones de entrenamiento con ciclo de estados & \texttt{sesiones\_entrenamiento} & Modelado \\
RF-19 & Asistencia por RFID o manual & \texttt{deportivo.asistencias} & Modelado \\
RF-20 & Evaluación diaria por criterios & \texttt{detalle\_evaluacion} & Modelado \\
RF-21 & Promedios de evaluación & \texttt{v\_promedio\_evaluacion} & Modelado \\
RF-22 & Notificación a representantes & --- & Planificado \\
\bottomrule
\end{longtable}
\renewcommand{\arraystretch}{1}

Los requisitos no funcionales (RNF-01 a RNF-13) se verifican con las
mediciones de la sección siguiente. La matriz
\ruta{docs/trazabilidad/matriz.csv} enlaza cada uno con su prueba y su
evidencia.

\subsection{Modelo de datos}

El esquema se divide en dos espacios de nombres con responsabilidades
separadas:

\begin{center}
\begin{tabular}{@{}llp{7.4cm}@{}}
\toprule
\textbf{Esquema} & \textbf{Tablas} & \textbf{Contenido} \\
\midrule
\texttt{seguridad} & 6 & Personas, usuarios, roles, relación usuario-rol,
estados y estudiantes. \\
\texttt{deportivo} & 9 & Posiciones, entrenadores, horarios, sesiones,
asistencias, criterios, evaluaciones, detalle y observaciones. \\
\bottomrule
\end{tabular}
\end{center}

Decisiones de modelado relevantes, todas materializadas como restricciones
en el motor y no solo como validación en la aplicación:

\begin{itemize}[leftmargin=*]
  \item \textbf{Baja lógica en lugar de borrado.} Las asistencias y
    evaluaciones referencian al estudiante por clave foránea; borrarlo
    físicamente destruiría su historial deportivo.
  \item \textbf{Unicidad parcial del código RFID.}
    \texttt{UNIQUE \ldots WHERE rfid\_codigo IS NOT NULL} permite muchos
    estudiantes sin credencial, pero impide que dos compartan la misma.
  \item \textbf{Una asistencia por sesión y estudiante.}
    \texttt{UNIQUE (id\_sesion, id\_estudiante)}.
  \item \textbf{Posición jugada por evaluación.}
    \texttt{detalle\_evaluacion} guarda la posición de \emph{ese} día,
    distinta de la habitual del estudiante, para que el histórico quede
    correctamente etiquetado cuando alguien juega fuera de su posición.
  \item \textbf{Coherencia temporal.}
    \texttt{CHECK (hora\_fin > hora\_inicio)} y
    \texttt{CHECK (dia\_semana BETWEEN 1 AND 7)} en los horarios.
\end{itemize}

El diccionario completo, con tipos, restricciones y disparadores, está en
\ruta{docs/basedatos/DATA-DICTIONARY.md}. El esquema no lo genera el ORM:
\texttt{ddl-auto} está fijado en \texttt{validate}, de modo que la aplicación
falla al arrancar si el esquema real y las entidades divergen.

% ============================================================
\section{Evidencia empírica}

Todas las mediciones de esta sección se generaron ejecutando el sistema real
levantado con \texttt{make up}, no en simulación, y sus artefactos crudos
están versionados.

\subsection{Rendimiento (Bloque C.1)}

Tres corridas independientes de k6 con 50 usuarios virtuales durante 30
segundos contra el \emph{endpoint} protegido \texttt{GET /api/estudiantes}:

\begin{center}
\begin{tabular}{@{}lrrrrr@{}}
\toprule
\textbf{Corrida} & \textbf{Media (ms)} & \textbf{Mediana} & \textbf{p90} & \textbf{p95} & \textbf{RPS} \\
\midrule
1 & 6,53 & 5,26 & 11,38 & 14,45 & 404,21 \\
2 & 6,67 & 5,94 & 11,63 & 15,32 & 403,52 \\
3 & 6,25 & 5,66 & 10,37 & 12,78 & 405,99 \\
\midrule
\textbf{Agregado} & \textbf{6,48} & --- & --- & \textbf{14,18} & \textbf{404,57} \\
\bottomrule
\end{tabular}
\end{center}

Con intervalo de confianza al 95\,\%: media $6{,}48 \pm 0{,}53$\,ms, p95
$14{,}18 \pm 3{,}20$\,ms, throughput $404{,}57 \pm 3{,}16$\,RPS, y
\textbf{0\,\% de errores} en las tres corridas.

El umbral exigido es p95 $<$ 200\,ms con caché caliente~\cite{iso25010}. El
resultado lo cumple con un margen de más de un orden de magnitud. \ok

\subsection{Seguridad: OWASP Top 10 (Bloque C.4)}

Seis controles verificados con peticiones reales contra el sistema en
ejecución~\cite{owasp2021}:

\begin{center}
\begin{tabular}{@{}llp{7.6cm}@{}}
\toprule
\textbf{Control} & \textbf{Resultado} & \textbf{Evidencia observada} \\
\midrule
A01 Control de acceso & \ok & \texttt{403} al operar sin el rol requerido \\
A02 Fallos criptográficos & \ok & TLS 1.3 en \texttt{:8443}; BCrypt coste 12 \\
A03 Inyección & \ok & \texttt{422} ante entrada malformada; sin SQL dinámico \\
A05 Configuración & \ok & CSP, HSTS, \texttt{nosniff}, \texttt{X-Frame-Options: DENY} \\
A07 Fallos de autenticación & \ok & Bloqueo exacto en el sexto intento (\texttt{429}) \\
A09 Registro y monitoreo & \ok & \texttt{AUTH\_LOGIN\_OK/FAIL} con IP, marca de tiempo y sujeto \\
\bottomrule
\end{tabular}
\end{center}

Los errores se devuelven como \texttt{ProblemDetail}~\cite{rfc9457}, sin
trazas de pila ni detalles internos. Ejemplo real capturado para A07:

\begin{lstlisting}[caption={Respuesta al sexto intento de autenticación}]
HTTP/1.1 429
Content-Type: application/problem+json
{
  "type":   "https://sged.uteq.edu.ec/errores/DemasiadasPeticiones",
  "title":  "Too Many Requests",
  "status": 429,
  "detail": "Demasiados intentos fallidos. Intente mas tarde."
}
\end{lstlisting}

\subsection{Cobertura de pruebas}

34 pruebas distribuidas en 7 clases, con \textbf{68,1\,\%} de cobertura de
instrucciones medida por JaCoCo (1073 cubiertas de 1575). El umbral exigido
es 60\,\%. \ok

\begin{center}
\begin{tabular}{@{}lrp{7.6cm}@{}}
\toprule
\textbf{Clase de prueba} & \textbf{Nº} & \textbf{Qué verifica} \\
\midrule
\texttt{AuthServiceTest} & 5 & Login correcto e incorrecto, registro, registro duplicado, disponibilidad. \\
\texttt{JwtServiceTest} & 4 & Emisor y audiencia válidos, audiencia distinta rechazada, token manipulado, refresco. \\
\texttt{LoginAttemptServiceTest} & 6 & Bloqueo al alcanzar el límite, no reinicio del TTL, borrado del contador al acertar. \\
\texttt{RedisBlacklistServiceTest} & 4 & Revocación con TTL restante, TTL no positivo ignorado, consulta de revocación. \\
\texttt{EstudianteControllerTest} & 9 & Los siete verbos del recurso más \texttt{404} y \texttt{422}. \\
\texttt{EstudianteServiceTest} & 5 & Paginación envuelta, baja lógica, persistencia transacional, delegación al procedimiento SQL. \\
\texttt{BackendApplicationTests} & 1 & Arranque del contexto contra el esquema migrado. \\
\midrule
\textbf{Total} & \textbf{34} & \\
\bottomrule
\end{tabular}
\end{center}

Un detalle relevante: durante esta entrega se descubrió que JaCoCo
\emph{nunca} había generado cobertura real, porque el \texttt{argLine}
configurado en Surefire sobrescribía el \emph{javaagent} del propio JaCoCo.
El informe anterior reportaba un valor que el sistema no estaba midiendo.

Las pruebas de \texttt{LoginAttemptServiceTest} merecen mención aparte:
verifican no solo que el bloqueo ocurra, sino que un fallo posterior
\emph{no} reinicie la ventana de tiempo. Sin esa segunda propiedad, un
atacante podría mantener la cuenta en estado de bloqueo indefinido, o bien
—según cómo se implemente— reiniciar el contador y evadir el límite.

\subsection{Calidad web: Lighthouse (Bloque C.5)}

Tres corridas con perfil móvil (Lighthouse v13.4.1):

\begin{center}
\begin{tabular}{@{}lrrl@{}}
\toprule
\textbf{Categoría} & \textbf{Media} & \textbf{Umbral} & \textbf{Estado} \\
\midrule
Rendimiento & 92,3 & $\geq$ 80 & \ok \\
Accesibilidad & 100 & $\geq$ 90 & \ok \\
Buenas prácticas & 96 & $\geq$ 90 & \ok \\
SEO & 63 & $\geq$ 90 & Ver §\ref{sec:seo} \\
\bottomrule
\end{tabular}
\end{center}

Métricas Web Vitals: LCP 2,86\,s; FCP 2,43\,s; TBT 25,7\,ms; y CLS
\textbf{0,000} en las tres corridas (la interfaz no desplaza contenido
durante la carga).

Esta medición expuso, y permitió corregir, tres defectos reales del
frontend: el documento declaraba \texttt{lang="en"} en una aplicación
íntegramente en español; conservaba el título \texttt{Frontend} generado por
Angular CLI y no tenía descripción; y carecía de un elemento \texttt{<main>},
lo que hacía fallar la auditoría de puntos de referencia de accesibilidad.
Corregidos los tres, la accesibilidad pasó de 97 a \textbf{100}.

Se corrigió además un defecto en la propia configuración de medición:
\ruta{lighthouserc.js} declaraba \texttt{preset: 'mobile'}, un valor que no
existe en Lighthouse, lo que abortaba toda corrida con código de salida 1.
Es decir, \textbf{esta medición nunca se había podido ejecutar}.

\subsubsection{Por qué el SEO de 63 es el resultado correcto}
\label{sec:seo}

El puntaje SEO bajó de 82 a 63 \emph{a causa de una corrección deliberada}.
Al añadir un \ruta{robots.txt} válido con \texttt{Disallow: /}, Lighthouse
activa la penalización \texttt{is-crawlable} («página bloqueada para
indexación»), que por sí sola descuenta 27 puntos, porque su categoría SEO
asume que el sitio \emph{quiere} ser indexado.

SGED es una aplicación de gestión interna que trata datos personales de
menores de edad. Que fuese indexable por buscadores sería un defecto de
privacidad, no una virtud. Elevar el SEO a 90 exigiría eliminar el
\ruta{robots.txt}: degradar la privacidad del sistema para mejorar una
métrica. Se optó por lo contrario, se relajó el umbral agregado de esa
categoría a advertencia con la justificación escrita en el propio archivo de
configuración, y se mantuvieron como bloqueantes las auditorías SEO que sí
aplican a una aplicación interna (\texttt{meta-description},
\texttt{document-title}, \texttt{html-has-lang}, \texttt{viewport}), todas
ellas en 1,0. \ok

\subsection{Usabilidad: SUS (Bloque C.3)}
\label{sec:sus}

\textbf{Medición pendiente.} El instrumento
(\ruta{docs/mediciones/sus/INSTRUMENTO-SUS.md}) contiene el cuestionario SUS
de diez ítems~\cite{brooke1996}, el protocolo con siete tareas trazadas a
requisitos concretos, y la escala adjetival de
interpretación~\cite{bangor2009}. El script \ruta{scripts/sus-analysis.py}
calcula media, desviación típica, intervalo de confianza al 95\,\%, mediana y
tasa de éxito por tarea.

El script se niega deliberadamente a ejecutarse con menos de diez
participantes. Las respuestas no han sido recolectadas todavía, y una muestra
fabricada invalidaría la medición. La aritmética del cálculo sí quedó
verificada contra los cuatro casos de referencia de la literatura (0, 50, 75
y 100 puntos).

% ============================================================
\section{Reproducibilidad}

El sistema se levanta completo desde una clonación limpia con un solo
comando:

\begin{lstlisting}[language=bash,caption={Arranque reproducible}]
git clone https://github.com/DarwinSM21/SGED_APPWEB.git
cd SGED_APPWEB
cp .env.example .env
make up
\end{lstlisting}

\begin{center}
\begin{tabular}{@{}ll@{}}
\toprule
\textbf{Objetivo} & \textbf{Acción} \\
\midrule
\texttt{make up} & Levanta el sistema completo \\
\texttt{make test} & JUnit 5 + cobertura JaCoCo \\
\texttt{make bench} & 3 corridas k6 + análisis con IC 95\,\% \\
\texttt{make audit} & Auditoría OWASP (6 controles) + SQL dinámico \\
\texttt{make clean} & Limpia contenedores, volúmenes y compilaciones \\
\bottomrule
\end{tabular}
\end{center}

Las imágenes de PostgreSQL y Redis están fijadas por digest SHA-256, no por
etiqueta móvil: dos construcciones del mismo \emph{commit} usan exactamente
los mismos binarios. La reproducibilidad se verificó clonando el repositorio
en un directorio independiente, no solo reiniciando el volumen local.

% ============================================================
\section{Consideraciones éticas}

El sistema trata datos personales de niños, niñas y adolescentes: nombres,
cédula, fecha de nacimiento, asistencia con hora y lugar, y evaluaciones
individuales de desempeño. Esa combinación es sensible porque los titulares
son menores que no pueden consentir por sí mismos, porque la asistencia es
información de localización temporal, y porque las evaluaciones constituyen
perfilado de personas~\cite{lopdp2021,acm2018}.

Se aplicó minimización: se decidió \emph{no} recolectar peso, altura,
contextura ni datos de alimentación, pese a haber sido considerados en el
diseño, por tratarse de datos de salud cuya recolección exigiría garantías
que este proyecto no puede sostener.

\ruta{docs/etica/ETHICS.md} documenta además cinco hallazgos abiertos, en
lugar de afirmar que todo está resuelto: la cédula se almacena sin validación
ni cifrado; las observaciones del entrenador son texto libre sin control
sobre un menor; no existe mecanismo de supresión de datos (la baja lógica
preserva el historial pero no es un borrado); el consentimiento del
representante no está modelado, lo que bloquea el módulo de notificaciones; y
el certificado TLS es autofirmado, adecuado para evaluación pero no para
producción.

Todos los datos del repositorio son ficticios. No se utilizaron datos reales
de menores en ninguna etapa del desarrollo ni de las mediciones.

% ============================================================
\section{Trazabilidad y honestidad académica}

La matriz \ruta{docs/trazabilidad/matriz.csv} enlaza cada requisito con su
historia de usuario, su caso de uso, su \emph{endpoint}, el archivo que lo
implementa, la prueba JUnit que lo cubre y la evidencia empírica
correspondiente.

El historial de \texttt{git} evidencia la autoría de cada integrante. Se
documenta en \ruta{CONTRIBUTORS.md} que dos \emph{commits} figuran bajo una
cuenta distinta por haber sido realizados desde un equipo de la universidad
con otra sesión configurada; la aclaración se hace por escrito en lugar de
reescribir el historial ya compartido. El mismo archivo declara
explícitamente el uso de herramientas de asistencia por inteligencia
artificial, siguiendo la guía de transparencia del SWEBOK
v4.0~\cite{swebok2024}.

% ============================================================
\section{Conclusiones}

\begin{enumerate}[leftmargin=*]
  \item El sistema cumple con holgura los umbrales cuantitativos exigidos:
    p95 de 14,18\,ms frente a un límite de 200\,ms, 68,1\,\% de cobertura
    frente a un mínimo de 60\,\%, y 6 de 6 controles OWASP verificados con
    evidencia reproducible.

  \item El aporte más sustantivo de esta entrega no fue añadir
    funcionalidad, sino \emph{verificar} lo que ya existía. Ese ejercicio
    reveló defectos que ninguna revisión de código había detectado: JaCoCo
    no medía cobertura en absoluto, los procedimientos almacenados estaban
    creados pero nunca se invocaban, la caché se rompía desde la segunda
    petición, y la configuración de Lighthouse impedía ejecutar la medición.
    Todos se corrigieron durante esta entrega.

  \item La distinción entre lo implementado y lo modelado se mantiene
    explícita en todo el documento. El dominio deportivo tiene su esquema
    migrado y sus restricciones de integridad definidas, pero no se presenta
    como funcionalidad entregada.

  \item Quedan abiertos dos frentes: la encuesta SUS, que requiere
    participantes reales, y la exposición REST del dominio deportivo, que es
    el objetivo natural de la Entrega Final.
\end{enumerate}

\newpage
\bibliographystyle{plain}
\bibliography{referencias}

\end{document}
